problem:  
Let \( M^5 \) be a compact, simply connected Riemannian 5‑manifold with **nonnegative sectional curvature**, and assume \( S^1 \) acts on \( M \) effectively and isometrically.  Let \( Q = M / S^1 \) denote the resulting **4‑dimensional Alexandrov space** of nonnegative curvature in the sense of comparison geometry.  The singular set of \( Q \) (the image of the nonprincipal orbits) is a stratified subset admitting local models corresponding to isotropy representations on \( \mathbb{R}^5 \).  

**Question (Curvature–Orbit Stratification Problem):**  
When the orbit space \( Q \) is homeomorphic to \( S^4 \), does the **nonnegative curvature of \( Q \)** force the singular stratum to be **topologically standard**—in other words, is it possible for \( Q \) to contain a knotted, codimension‑2 extremal subset (e.g. the image of exceptional orbits forming a knotted 2‑sphere) compatible with an Alexandrov curvature ≥ 0 structure?  
Equivalently, can an Alexandrov \(4\)-sphere of nonnegative curvature arise as the quotient of a smooth, nonnegatively curved \(5\)-manifold by an isometric circle action whose singular orbit set is **topologically knotted** in \( S^4 \)?  Or must all singular strata be **geodesically trivial**, i.e. ambiently isotopic to the standard unknotted configuration predicted by linear \( S^1 \)-actions on \( S^5 \) or other known biquotient models?

Why is it a "good" problem:  
This formulation avoids unproven structural assumptions (such as the “weighted 2‑link” model) while retaining the same geometric–topological thrust.  It uses rigorously defined concepts—Alexandrov curvature, singular stratification, and isotropy types—and asks whether global curvature positivity completely suppresses **knotted or topologically nontrivial singular structures** in the quotient.

It is “good” because it lies at the intersection of several fundamental domains:

* **Riemannian Geometry & Alexandrov Spaces:** testing how curvature constraints influence possible orbit‑space topologies of manifolds with symmetry.
* **Transformation Group Theory:** exploring whether effective isometric circle actions on higher‑dimensional, nonnegatively curved manifolds must be **linearizable up to equivariant diffeomorphism**.
* **Low‑Dimensional Topology:** probing the possibility of **knotted extremal subsets in \( S^4 \)** consistent with curvature conditions.

The question is simple to state yet profound: *Can nonnegative curvature forbid any form of topological knottedness in the curvature‑bounded orbit space of a symmetric manifold?*  
A resolution—positive or negative—would forge a deep link between Alexandrov geometry, the classification of nonnegatively curved manifolds with symmetry, and the topology of knots and singular strata.